\title{Large Language Models Can Automatically Engineer Features for Few-Shot Tabular Learning}

\begin{abstract}
Large Language Models (LLMs) possess impressive capabilities for addressing complex and previously unencountered reasoning tasks, making them highly promising for tabular learning—a cornerstone in numerous practical domains. In this work, we introduce a new in-context learning paradigm named \textsf{FeatLLM}, which leverages LLMs to engineer features, thereby creating a dataset tailored for improved performance on tabular prediction problems. The features derived by \textsf{FeatLLM} are subsequently utilized with straightforward downstream predictors, such as linear regression models, achieving notable results in few-shot learning scenarios. Importantly, the \textsf{FeatLLM} approach confines inference solely to these simple predictive models combined with the features synthesized by the LLMs, eliminating any requirement to interact with the LLM for each inference instance. Unlike other LLM-based methods, \textsf{FeatLLM} circumvents prompt size restrictions and requires only API-level access to LLMs. Extensive experimentation on diverse tabular datasets illustrates that \textsf{FeatLLM} is capable of generating superior feature representations and decision rules, consistently outperforming other state-of-the-art techniques like TabLLM and STUNT by an average margin of 10\%.
\end{abstract}

\section{Introduction}
Large Language Models (LLMs) have emerged as powerful tools for generating responses to novel tasks, often excelling even in the absence of task-specific fine-tuning~\cite{creswell2022selection,imani2023mathprompter,taylor2022galactica}. By virtue of training on massive corpora of text, these models acquire substantial background knowledge, enabling them to act as substitutes for specialized human expertise across numerous fields~\cite{Genomes2021,singhal2023large,wilcox2003role}, thus fostering automation in applications as varied as conversational analytics~\cite{finch2023leveraging} and automatic program repair~\cite{fan2023automated}. A notable aspect of LLMs is their proficiency at interpreting prompts that combine task instructions with a handful of illustrative examples; through such prompt engineering, these models can discern context and identify salient variables or samples~\cite{brown2020language,zhao2023survey}. This characteristic positions LLMs as effective automated feature engineers capable of insightful data analysis.

Building on these strengths, LLMs' inherent knowledge and inferential abilities have been increasingly leveraged in tabular data learning settings. For example, domain-specific insights—such as recognizing a correlation between advanced age and heightened disease risk—prove valuable for applications like disease prediction. Approaches in this space typically articulate tasks and features in natural language, transform tabular records into serialized sequences, and supply them as input to an LLM~\cite{dinh2022lift,wang2023anypredict}. These workflows may entail in-context learning strategies~\cite{nam2023semi} or utilize lightweight fine-tuning methodologies~\cite{dinh2022lift,hegselmann2023tabllm}. The infusion of LLM-supplied priors has delivered noteworthy improvements over classical tabular methods, particularly under limited-data conditions~\cite{hegselmann2023tabllm}.

Existing tabular learning approaches utilizing LLMs face notable challenges. Achieving end-to-end predictions typically demands running the LLM inference separately for every individual sample, resulting in substantial computational overhead. Furthermore, most of these approaches depend on fine-tuning the LLM to reach high prediction accuracy. This reliance restricts their applicability to scenarios where the LLM's training interfaces or fine-tuning APIs are not publicly accessible—a situation made more pressing as leading LLMs today tend to only offer limited inference API access~\cite{anil2023palm,openai2023gpt4}. Another critical limitation arises when prompts become excessively long~\cite{liu2023lost}: as the number of tabular features grows, the input text length can render these methods impractical or cause a notable drop in performance. Collectively, these issues limit the feasibility of deploying LLM-based models for real-world tabular data applications.

To address these challenges, our method reimagines the function of LLMs within tabular prediction workflows. Rather than relying on LLMs for direct prediction, we harness their capabilities for feature engineering—utilizing their extensive background knowledge in a more computationally efficient manner. We present a new in-context learning paradigm, \textsf{FeatLLM}, designed to extract the underlying "criteria" from LLM decision-making. For example, in clinical prediction tasks, an LLM can articulate decision rules specifying feature thresholds associated with specific diseases. These derived criteria serve as the basis for constructing new features that can supplant the original variables for downstream learning tasks. This methodology shifts the LLM’s value-add from end-to-end modeling to informative feature creation, allowing simpler predictive models to replace complex LLM-based inference and consequently improving inference speed. By leveraging in-context learning—embedding demonstrations into prompts—feature extraction can be performed without any model fine-tuning. To overcome issues stemming from prompt length, we incorporate ensemble-inspired techniques such as feature bagging~\cite{breiman2001random}, effectively reducing the information packed into any single prompt.

\textsf{FeatLLM} organizes the process of prompt-based feature engineering into two sequential components: (1) comprehending the task at hand and identifying how features relate to the prediction targets; and (2) utilizing insights from the initial stage, along with a limited set of annotated examples, to derive discriminative rules for class separation. This step-wise reasoning method guides LLMs in disregarding irrelevant features during rule formation. The resulting rules are systematically applied to the dataset (interpreting them as executable code), which subsequently encodes each sample with a binary indicator reflecting compliance with the corresponding rule. A simple model is then trained on these engineered features to approximate class membership probabilities. The overall stability is enhanced by adopting an ensemble methodology, which entails repeated feature construction together with feature bagging. By tuning the quantities of features and samples selected during the bagging process, the difficulty posed by large prompt sizes can be alleviated. Notably, \textsf{FeatLLM} operates without any additional model training and incurs minimal inference overhead, thereby broadening the scope of LLM applicability to a variety of practical scenarios.

To assess its efficacy, we benchmark \textsf{FeatLLM} on 13 tabular datasets under low-data conditions, where it consistently demonstrates high and reliable performance. Our findings indicate that LLMs effectively integrate prior knowledge with dataset-specific information, enabling the extraction of valuable features. Our method surpasses several state-of-the-art few-shot learning approaches in multiple experimental settings. The implementation is publicly available through an anonymized GitHub repository at~\texttt{https://github.com/Sungwon-Han/FeatLLM}.

\section{Related Work}

\subsection{Few-Shot Learning with Tabular Data} 

Few-shot learning has made it possible to develop representations that generalize well, even when only limited labeled data are available~\cite{chen2018closer}. While research in this area originally centered around image data, the underlying principles have been successfully adapted to multiple other domains, such as text, audio, and more recently, tabular datasets~\cite{majumder2022few,nam2023stunt,schick2021exploiting}. Nevertheless, learning from a handful of labeled examples is prone to the pitfall of capturing spurious patterns that do not generalize~\cite{bartlett2021deep}. To mitigate these issues, contemporary approaches have leveraged supplementary data sources or incorporated domain-aware prior knowledge, thereby instilling suitable inductive biases during model training. Techniques include harnessing copious amounts of unlabeled tabular data in conjunction with tailored data augmentations to enable semi-supervised learning~\cite{ucar2021subtab,yoon2020vime}, or obtaining advantageous model initializations via task-independent unsupervised pre-training schemes~\cite{somepalli2021saint}. Such unsupervised pre-training methods often utilize strategies like data masking or corruption paired with a reconstruction loss~\cite{arik2021tabnet,majmundar2022met}, or integrate data augmentation to facilitate contrastive learning objectives~\cite{bahri2022scarf,chen2020simple}. However, because tabular datasets exhibit high heterogeneity across features and distributions, designing universal augmentations is challenging, resulting in limited efficacy of these techniques in few-shot regimes~\cite{nam2023stunt}. 

More recently, approaches have shifted toward pre-training models on a heterogeneous collection of tabular datasets—often unrelated to the downstream task—and subsequently fine-tuning or adapting them to the target application via transfer learning~\cite{zhang2023generative,levin2022transfer,zhu2023xtab}. As an illustration, TransTab~\cite{wang2022transtab} utilizes transformer architectures to learn semantic relationships across columns, leveraging not only the cell values but also column headers to capture cross-table knowledge. Such cross-table representations have empowered models to perform inference on new tabular problems in both zero-shot and few-shot settings. In another example, TabPFN~\cite{hollmann2022tabpfn} pre-trains a Bayesian neural network using mixtures of artificial distributions, generating synthetic tabular data. This model then conducts inference for the target task by evaluating training and test examples together, requiring no additional parameter updates. Prior knowledge acquired from varied tabular distributions allows these methods to outperform conventional tree-based models, particularly when annotated data is scarce.

\subsection{Language-Interfaced Tabular Learning} 

The integration of large language models (LLMs)~\cite{anil2023palm,openai2023gpt4} represents another promising route for introducing prior knowledge into tabular learning tasks. LLMs, having been exposed to vast and varied textual data, demonstrate the capacity to generalize across novel tasks and domains~\cite{brown2020language}. Recent work explores leveraging the implicit knowledge contained within LLMs by converting tabular data into textual form and presenting this as input prompts~\cite{dinh2022lift,hegselmann2023tabllm,wang2023anypredict}. These prompts typically consist of serialized table values, feature metadata, and task descriptions, allowing LLMs to reason about associations between task objectives and features in order to make predictions. This research trajectory has bifurcated into methods that either fine-tune language models using parameter-efficient strategies such as LoRA~\cite{hu2021lora} or IA3~\cite{liu2022few}, or employ in-context learning approaches, where task demonstrations and a few annotated examples are directly incorporated into the prompt—bypassing additional training~\cite{dinh2022lift,hegselmann2023tabllm,nam2023semi,slack2023tablet}.

While prior approaches have shown considerable success in enhancing few-shot tabular learning, their reliance on routing inference through the LLM for every instance introduces limitations for real-world deployment in industrial settings. This work shifts away from the standard end-to-end, black-box paradigm where each individual sample is processed by an LLM for prediction. Instead, the focus is on harnessing the LLM to infer interpretable rules or conditions corresponding to each class directly. \textsf{FeatLLM} translates these generated rules into programmatic features, allowing for predictions to be made without further interaction with the LLM. In this manner, in-context learning is employed to derive rules by drawing upon both the model’s prior knowledge and information provided by the few-shot exemplars.

\section{Methods}
The goal of \textsf{FeatLLM} is to derive a collection of ``rules"—the defining criteria associated with each possible class—from a limited number of provided samples, as illustrated in Figure~\ref{fig:main_model}. These rules, generated using an LLM, are subsequently parsed and converted into new binary features, which indicate whether each rule holds for a particular input sample. These binary features are then aggregated through a dot product with non-negative, learnable parameters in order to approximate the probability of membership in each class. Through this training process, the model learns to weigh the influence of each extracted feature in contributing to class predictions (see Section~\ref{Sec:training}). To enhance robustness and predictive power, the procedure is iteratively performed in a bagging framework and the results are ensembled. The precise formulation of the problem and details for each component of the methodology are outlined in the following sections.

\vspace*{-0.12in}
\paragraph{Problem formulation.} We begin with a tabular dataset $\mathcal{D} = \{(\mathbf{x}^i, \mathbf{y}^i)\}_{i=1}^{N}$, which contains $N$ labeled data points. Each data point comprises $d$ input features, where each feature vector $\mathbf{x}^i$ is $d$-dimensional. The dataset also includes natural language descriptions for each feature, noted as $F = \{f_j\}_{j=1}^{d}$, with corresponding definitions listed separately. Within the context of classification, each label $\mathbf{y}^i$ belongs to a finite set of possible classes. In $k$-shot learning settings, we randomly select a subset of $k$ samples, where $k<N$, to use for model training.

\subsection{Prompt Design for Extracting Rules}
\label{Sec:prompt_design}
To support rule extraction from large language models (LLMs) through more sophisticated reasoning, we craft the solution pathway to reflect the analytical strategies an expert would apply when engaging with tabular data. The prompt used for this purpose contains three primary sections (see Fig.~\ref{fig:prompt_for_rule} and Appendix~\ref{sec:example_rule}):

\vspace*{-0.12in}
\paragraph{Basic information description.} This segment conveys critical context needed for resolving the task (highlighted in orange in Fig.~\ref{fig:prompt_for_rule}). It specifies the problem at hand along with class label details, provides descriptions for each input feature, and includes illustrative examples. The objective of the task is expressed as a query—such as: ``Does this patient's myocardial infarction complications data indicate chronic heart failure? Yes or no?"—with additional examples detailed in Appendix~\ref{sec:task_desc}. Feature descriptions state their respective data types and may contain definitions; for categorical features, representative values may also be listed. To serve as demonstrations, a small set of annotated training examples is converted into textual form along with their true labels. The serialization method mirrors that of prior work~\cite{dinh2022lift,hegselmann2023tabllm}:

\begin{align}
&\text{Serialize}(\mathbf{x}^i,\mathbf{y}^i, F) = \nonumber \\ 
&``\ f_1\ \text{is}\ \mathbf{x}^i_1.\ ... \ f_d\  \text{is}\  \mathbf{x}^i_d.\ \ \text{Answer:}\  \mathbf{y}^i\ ”
\end{align}

\vspace*{-0.12in}
\paragraph{Reasoning instruction.} Our objective is to strengthen the LLM's reasoning capabilities by supplying explicit reasoning instructions (illustrated as blue text in Fig.~\ref{fig:prompt_for_rule}). The instruction begins with a prompt akin to chain-of-thought prompting~\cite{wei2022chain}, setting the stage for subsequent reasoning. The reasoning process is systematically separated into two stages: Initially, the LLM is prompted to infer possible causal relationships or tendencies between the features and the target task, utilizing its background knowledge and common sense—deliberately without referencing any concrete example demonstrations. This ensures the LLM's conclusions in this stage draw upon prior general understanding of the task. Next, in the second stage, the LLM synthesizes the information from the first stage along with the given example demonstrations to devise class-specific rules. Such a bifurcated approach is intended to minimize the likelihood of the LLM detecting superficial or irrelevant feature correlations, thereby promoting focus on the most impactful features. To maintain consistency and prevent issues such as underfitting (too few rules) or unnecessarily verbose prompts (too many rules), we specify a fixed count of rules per class (i.e., 10)\footnote{A discussion on how the number of rules affects performance is available in Section~\ref{sec:analysis}.}. Furthermore, the LLM is directed to reference the feature types provided in the basic information section while formulating rules, to mitigate mismatches in feature type assignments.

\vspace*{-0.12in}
\paragraph{Response instruction.} For clarity and ease of subsequent rule interpretation, we guide the LLM to format its outputs according to tailored response instructions (depicted as yellow text in Fig.~\ref{fig:prompt_for_rule}). The prompt specifies both the overall answer structure expected for each class (i.e., the prescribed response format) and the formatting conventions for each individual rule (i.e., the format for the rule). This framework enables the LLM to appropriately adjust its outputs based on the respective feature types. Additionally, without explicit instruction, the LLM may integrate several rules via logical connectors such as AND and OR. An illustrative result of this structured response approach is given in Appendix~\ref{sec:outcome-rule}.

\subsection{Inferring Class Likelihood via Rules}
\label{Sec:training}

\paragraph{Parsing rules for feature generation.} The rules derived in the earlier phase serve as the basis for constructing new binary features. Specifically, for every class, these features take a value of ‘0’ or ‘1’, reflecting whether a sample conforms to the respective class-specific rules. Such features are then utilized to assess the likelihood of a sample’s class membership. Since the LLM produces these rules in natural language form, it is essential to design a parsing mechanism that can automate feature extraction. Although rule formatting is controlled during generation to aid parsing, the LLM output often contains minor syntactic inconsistencies or noise~\cite{kovavc2023large}. For example, the LLM may replace mathematical symbols such as ``$>$'' and ``$<$'' with verbal equivalents like ``greater than'' or ``less than,'' complicating parsing processes. 

To overcome the difficulties of handling noisy textual rules, rather than implementing intricate parsing logic manually, we instead utilize the LLM’s own capabilities. The LLM excels at capturing semantic meaning and is also adept at producing executable code when given suitable prompts, as it is trained on code data. We thus provide it with the function signature, as well as descriptions of inputs, outputs, and the inferred rules, packaging this information into a prompt that is fed into the LLM. Examples of such prompts and the corresponding outputs can be seen in Figures~\ref{fig:prompt_for_parsing} and~\ref{sec:example_parsing}-\ref{sec:example_parse_outcome} in the Appendix. The resulting code is run through Python’s \texttt{exec()} function to complete the data transformation step.

\vspace*{-0.12in}
\paragraph{Inferring class likelihood.} Once the rule-derived binary features are obtained for each class, one straightforward approach to estimate the likelihood that a sample belongs to a given class is by tallying the number of satisfied rules per class—essentially, summing up the class-specific binary features. Nevertheless, individual rules may vary in their predictive strength, making it important to learn the relative contribution of each feature using data. To accomplish this, we employ a standard machine learning (ML) model, training it on each class’s set of binary features to determine their importances. As an illustration, one might utilize a bias-free linear model for this purpose. Denote the binary feature vector for sample $\mathbf{x}^i$ and class $k$ as $\mathbf{z}_k^i \in \{0, 1\}^R$, where $R$ represents the count of rules per class and $c$ is the number of classes in total. The class probabilities $\mathbf{p}^i$ are then evaluated as follows:

\begin{align}
\text{logit}_k^i &= \max(\mathbf{w}_k, 0) \cdot \mathbf{z}_k^i, \nonumber \\
\mathbf{p}^i &= \text{Softmax}({[\text{logit}_{1}^i, ..., \text{logit}_{c}^i]}). \label{eq:infer_prob}
\end{align}

In this context, $\mathbf{w}_k$ denotes the adjustable parameter within the linear model, encapsulating the relevance assigned to each individual feature. Restricting these weights to the positive orthant guarantees that features extracted from the LLM cannot diminish the probability assigned to any given class during training. Subsequently, the resulting logit vector is mapped to class probability distributions via the softmax operation. To estimate $\mathbf{w}_k$, we employ cross-entropy loss, fitting the weights to a limited set of labeled examples in a few-shot setting.

\vspace*{-0.12in}
\paragraph{Ensembling with bagging.} To bolster predictive performance, we replicate the described process multiple times, thereby generating a suite of inference models whose outputs are combined into an ensemble decision. For $T$ independent training instances yielding weight vectors $\{\mathbf{w}[t]\}_{t=1}^T$, final class probability predictions are obtained by taking the mean across the individual probability vectors $\mathbf{p}^i$, each evaluated using a different $\mathbf{w}[t]$, as articulated in Eq.~\ref{eq:infer_prob}.

Diversity within the ensemble is promoted by setting a relatively high temperature parameter during LLM inference, which introduces variation into the rule generation process. Moreover, the sequence of few-shot demonstration examples is shuffled for each trial, leveraging evidence that order effects in prompt design can influence LLM outputs~\cite{lu2022fantastically}\footnote{Further discussion of rule diversity is provided in Appendix~\ref{sec:diversity}}. To further exploit the heterogeneity among individual models for greater robustness, we apply bagging: in each run, only a subset of features or samples is selected. This mechanism proves especially beneficial when confronted with a large pool of training instances or numerous features. The proposed ensembling strategy confers two notable strengths. First, the ensemble can absorb inaccuracies stemming from any mis-specified rules in some individual trials by leveraging correct rules from others, thereby increasing system resilience—an effect related to the self-consistency property of LLMs~\cite{wang2022self}, which suggests that frequently emergent rules across different runs are more trustworthy. Second, this ensembling framework helps mitigate issues stemming from LLM prompt length limitations, as bagging reduces input size and thus maintains system stability even with expansive tabular datasets. While any single ensemble constituent may overlook particular feature relationships, aggregating multiple weak learners drawn from distinct trials allows the ensemble as a whole to compensate for incomplete coverage or underrepresentation at the individual model level.

\section{Experiments}
We assess the performance of \textsf{FeatLLM} across a variety of tabular datasets and perform a series of analyses to understand the inner workings of our framework. Due to limitations on space, supplementary experimental results—such as alternative LLM configurations, qualitative studies, and more—are presented in the Appendix.

\subsection{Performance Evaluation}
\paragraph{Datasets.}
Our experimental setup includes 13 datasets covering binary and multi-class classification: (1) Adult~\cite{asuncion2007uci}, which involves predicting if an individual’s annual income exceeds \$50,000; (2) Bank~\cite{moro2014data}, focused on determining whether a client will opt for a term deposit; (3) Blood~\cite{yeh2009knowledge}, concerning prediction of repeat blood donation; (4) Car~\cite{kadra2021well}, assessing the overall quality classification of vehicles; (5) Communities~\cite{misc_communities_and_crime_183}, tasked with forecasting crime rates within specific regions; (6) Credit-g~\cite{kadra2021well}, assessing creditworthiness; (7) Diabetes\footnote{\texttt{kaggle.com/datasets/uciml/pima-indians-diabetes-database}}, to diagnose diabetes occurrence; (8) Heart\footnote{\texttt{kaggle.com/datasets/fedesoriano/heart-failure-prediction}}, for predicting coronary artery disease; and (9) Myocardial~\cite{misc_myocardial_infarction_complications_579}, predicting the incidence of chronic heart failure.

Moreover, we incorporate both newly published datasets and synthetic datasets—less commonly featured in literature and excluded from LLMs’ pretraining corpus: (10) Cultivars, for classifying soybean varieties by grain yield\footnote{\texttt{archive.ics.uci.edu/dataset/913}}; (11) NHANES, for predicting an individual’s age group from provided records\footnote{\texttt{archive.ics.uci.edu/dataset/887}}; (12) Sequence-type, a four-class problem involving sequence type identification (arithmetic, geometric, Fibonacci, Collatz); and (13) Solution-mix, which involves predicting whether the mixture concentration from four solutions surpasses a 0.5 threshold. Cultivars and NHANES were only recently added to the UCI Machine Learning Repository after September 2023, ensuring their absence from the pretraining data of the LLM API (gpt-3.5-turbo-0613) utilized in our work. The Sequence-type and Solution-mix datasets are synthetically generated, guaranteeing they could not be present in the LLM API’s memory, thus validating our experimental evaluation.

The datasets employed exhibit a wide range in terms of size and intricacy, as summarized in Table~\ref{Tab:basic-desc}. Each dataset provides unambiguous attribute names and accompanying descriptions. Notably, ‘Communities’ and ‘Myocardial’ possess over 100 attributes, creating context window challenges for the LLM models.

\vspace*{-0.12in}
\paragraph{Baselines.}
We benchmark our method against nine different baselines. The first three represent standard methods for tabular data analysis: (1) Logistic Regression (LogReg); (2) XGBoost~\cite{chen2016xgboost}; and (3) RandomForest~\cite{ho1995random}. Two subsequent baselines exploit access to unlabeled data: (4) SCARF~\cite{bahri2022scarf}; and (5) STUNT~\cite{nam2023stunt}. It should be noted that gathering ample unlabeled data can be practically difficult. The next, (6) TabPFN~\cite{hollmann2022tabpfn}, uses synthetically generated data of varying distributions for pretraining. The final three baselines are grounded in LLM usage: (7) In-context learning (In-context)~\cite{wei2022emergent}, where a handful of training examples are directly inserted into the input prompt; (8) TABLET~\cite{slack2023tablet}, which augments prompts with additional rule-based and prototype-derived cues via an external classifier to improve prediction; and (9) TabLLM~\cite{hegselmann2023tabllm}, leveraging parameter-efficient tuning approaches such as IA3~\cite{liu2022few} to enable LLM adaptation with limited annotated samples.

\vspace*{-0.12in}
\paragraph{Implementation details.}
\textsf{FeatLLM} is built upon the GPT-3.5 architecture as its LLM core, though the framework is fundamentally adaptable to alternative LLM backbones (refer to Appendix~\ref{sec:backbones} for comparative results across different backbones). For LLM inference, we apply a temperature setting of 0.5, and retain the API's default top-p value of 1. In our experiments, we configure the ensemble size to 20 and set the rule extraction count to 10. Further analysis of how hyper-parameter choices affect performance is presented in Figure~\ref{fig:hyperparameter2} and detailed in Appendix~\ref{sec:hyper-parameter}.

The linear model leverages the Adam optimizer with a learning rate of 0.01 and is trained over 200 epochs. Model selection is conducted through $k$-fold cross-validation to identify the best epoch. When evaluating baselines, we use GPT-3.5 for in-context learning experiments, and employ T0~\cite{sanh2022multitask} within TabLLM as per its original configuration. It is important to highlight that TabLLM limits LLM usage to models with open-source checkpoints, excluding the use of GPT-3.5. For conventional machine learning approaches such as LogReg, XGBoost, and RandomForest, their hyper-parameters are optimized via Grid-search in combination with $k$-fold cross validation. The value of $k$ is selected as either 2 or 4 to ensure that every class is represented in the training data. Comprehensive implementation details can be found in Appendix~\ref{sec:baselines}.

\vspace*{-0.12in}
\paragraph{Results.}
Tables~\ref{tab:main1} and \ref{tab:main2} display comparative performance metrics for a variety of datasets. Each experiment was conducted three times, varying the random splits for training and testing, and the mean along with the standard deviation of the AUC scores are reported. Across the board, \textsf{FeatLLM} either leads or closely trails the leading model among the baseline approaches evaluated. Remarkably, our method delivers performance on par with traditional tabular learning baselines using only 4 shots, which is equivalent to what those baselines achieve with 32 shots (see Fig.~\ref{fig:compares}). While STUNT and TabPFN yielded notable gains on particular datasets, their cumulative advantage over classic machine learning baselines was limited. Furthermore, LLM-based models, such as those relying on in-context learning, TABLET, or TabLLM, were unable to process tabular datasets characterized by high dimensionality. The integration of bagging and ensemble strategies in our framework demonstrated reliable efficacy, maintaining robust performance even when the dataset involved a large number of features. It is worth mentioning that our approach to bagging and ensemble methods can potentially be extended to other LLM-centric models; however, this would entail a twofold increase in per-sample inference cost, markedly escalating computational demands and rendering such an extension largely infeasible. 

\subsection{Analysis \& Discussion}
\label{sec:analysis}
\paragraph{Ablation study.} We investigate the influence of critical components through ablation studies, considering: (1) \textsf{-Tuning}: excluding the linear model’s weight tuning step and instead aggregating binary features directly to compute logits; (2) \textsf{-Ensemble}: removing the ensemble mechanism; (3) \textsf{-Description}: omitting feature descriptions; and (4) \textsf{-Reasoning}: eliminating Step-1 from the reasoning instruction during rule generation.

Table~\ref{tab:ablation} provides the mean change in AUC resulting from each ablation, averaged over all datasets. In every instance, performance declined when a component was ablated. Notably, the importance of weight tuning becomes more pronounced as the number of shots increases, indicating that, with ample data, fine-tuning rule weights can significantly enhance performance. In contrast, when the shot number is low, omitting feature descriptions and reasoning steps yields a larger negative impact, implying that leveraging the LLM’s inherent prior knowledge is particularly important in data-scarce settings. Finally, our ensemble method shows consistent performance improvement in every evaluated scenario.

\vspace*{-0.12in}
\paragraph{Training \& inference time comparison.} 
We evaluate and contrast the training and inference durations among various approaches. The assessment is conducted using the Adult dataset, utilizing a training subset containing 16 examples. Unless otherwise stated, a single A100 GPU is utilized for computation, except in the case of TabLLM, which leverages four A100 GPUs to enable model parallelism. Table~\ref{tab:cost} displays a summary of the runtime outcomes. Importantly, \textsf{FeatLLM} achieves a notably low inference latency, on par with traditional techniques. By contrast, alternative LLM-driven strategies, including In-context learning, TABLET, and TabLLM, tend to result in substantially greater inference times. Regarding training overhead, \textsf{FeatLLM} is similarly efficient as other LLM-oriented models; it makes only 30 API requests, which is modest in terms of cost and can also be executed in parallel, thereby minimizing budgetary concerns.

\vspace*{-0.12in}
\paragraph{Quality of features generated by \textsf{FeatLLM}.} 
To assess the informativeness of the rule-derived features created by \textsf{FeatLLM} for real-world applications, we conduct a targeted experiment. In particular, we compute the mean AUROC between these newly derived features and the outcome variable, and subsequently determine the proportion of features for each dataset that surpass an AUROC of 0.5. 
Table~\ref{tab:quality} presents these findings. The analysis demonstrates that most constructed rules meaningfully correspond to the target attribute and contribute valuable information to the downstream task (see the ``Before tuning'' column). Additionally, during the linear model tuning process, rules deemed less pertinent (those assigned a weight beneath a predefined threshold) are systematically discarded (refer to the ``After tuning'' column). The threshold for pruning was chosen as 0.1, reflecting an analysis of the global distribution of learned weights.

\vspace*{-0.12in}
\paragraph{Effect of adding spurious correlations.} We conducted an investigation to evaluate how well different approaches can disregard irrelevant, noisy correlations under low-shot conditions. Specifically, we selected the Heart dataset as our testbed and supplemented it with randomly chosen columns from the Adult dataset—columns that bear no relevance to the predictive task associated with the Heart dataset. The performance of various models was then assessed as we systematically increased the quantity of these irrelevant columns. To ensure a fair comparison, each method was provided with 8 labeled examples, and we did not incorporate any unlabeled data. As a result, approaches that require such data (e.g., SCARF or STUNT) were omitted from this analysis.

Figure~\ref{fig:spurious} illustrates how model performance shifts in response to the introduction of spurious correlations. Among the methods tested, \textsf{FeatLLM} exhibited the smallest reduction in accuracy as more spurious correlations were introduced. This resilience appears to stem from the framework's design, which relies on leveraging prior knowledge to relate features and tasks when deriving rules. As such, rules constructed from irrelevant features are largely avoided, enhancing robustness. Quantitatively, we noted that rules using noisy features comprised merely 1-2\% of the overall rule set. Nevertheless, it is also evident that \textsf{FeatLLM} is not immune to learning misleading associations: if the added columns are from domains closely linked to the Heart dataset, the model may still conflate spurious with causal relationships (refer to Appendix~\ref{sec:spurious-appendix}).

\vspace*{-0.12in}
\paragraph{Hyper-parameter analysis}
\textsf{FeatLLM} is primarily governed by two hyper-parameters: the ensemble count and the number of rules retrieved per class during each LLM interaction. We systematically assessed how these parameters influence the model’s effectiveness. The experimental findings reveal that increasing the ensemble size consistently enhances model accuracy, though performance gains plateau once the number exceeds approximately 20 (refer to Fig.~\ref{fig:appendix-hyperparameter1} in Appendix). Regarding the rule count per class, although we did not observe a statistically meaningful performance variation, a selection of 10 rules was identified as optimal, balancing between the length of the prompt and the resulting accuracy (see Fig.~\ref{fig:hyperparameter2}). We hypothesize that while incorporating a greater number of rules supplies richer information due to increased input dimensionality, it may also promote overfitting, particularly under the low-shot scenario.

\section{Conclusion}
We introduce \textsf{FeatLLM}, which marks a substantial advancement in leveraging LLMs for few-shot learning on tabular datasets. Unlike existing methods that simply rely on LLMs for prediction at the instance level, \textsf{FeatLLM} is centered on constructing prediction guidelines by functioning as a feature engineering framework. Through the processes of rule induction and bagged ensemble modeling, \textsf{FeatLLM} achieves competitive accuracy with minimal inference overhead, operates using API access alone without necessitating additional model training, and effectively handles datasets with larger feature spaces. This framework provides a robust and resource-efficient strategy for applying LLMs to tabular data and consistently delivers strong performance on practical datasets.

Currently, \textsf{FeatLLM} targets the low-shot learning regime. Going forward, our objectives include expanding its capacity to synthesize features for datasets containing greater numbers of instances, incorporating more diverse types of feature representations beyond rule-based formulations, and enhancing the interpretability of its outputs, with the goal of facilitating its adaptation to broader problem domains.

\section{Broader Impact}
The presented approach, \textsf{FeatLLM}, is designed to exploit the domain knowledge encapsulated within LLMs for improving the effectiveness of tabular learning tasks, surpassing what conventional supervised learning can achieve in data-sparse environments. By reinforcing data efficiency through the integration of prior knowledge, this work contributes substantially to alleviating the overfitting challenges often encountered with small datasets. The methodology is especially advantageous for sectors such as finance and healthcare, where data annotation is labor-intensive or otherwise difficult, and where pretrained LLMs can bring in considerable background knowledge drawn from domain-specific corpora. For broader deployment, it will be important to rigorously investigate the influence of societal biases or inaccurate information present in the LLM’s pretrained knowledge, as this prior information could substantially impact—or even outweigh—the influence of newly observed labeled samples in prediction outcomes.

\end{document}